\chapter{Implementation and Testing}
\label{chap:implementation}

This chapter provides a comprehensive technical account of the algorithmic implementations
underpinning Vital Guardian, the verification strategy for each sub-module, and the
extensive ablation studies that guided architectural decisions throughout the project.

% =============================================================
\section{Core Algorithmic Implementations}
\label{sec:algorithms}

The production system is composed of three tightly integrated tiers:
\textbf{Edge Vision Processing}, \textbf{Edge Audio Processing}, and
\textbf{LLM-Driven Cognitive Verification}. Each tier is described in detail below.

% -------------------------------------------------------------
\subsection{Visual Guardian: YOLOv11n Person Detection}
\label{subsec:yolo}

Person detection is the first step in the vision pipeline. A YOLOv11n model,
pre-compiled and optimized via Intel OpenVINO, scans every incoming frame and
returns a bounding box around the patient. The cropped region of interest (ROI)
is then forwarded to the downstream classifiers, eliminating background noise and
reducing computational load.

To maintain real-time throughput on CPU hardware, a frame-skipping strategy is
applied: the detector runs once every $N$ frames (configurable via
\texttt{PERSON\_DETECTOR\_PROCESS\_EVERY}), caching the last valid bounding box
for intermediate frames.

\begin{algorithm}[H]
\caption{YOLOv11n-OpenVINO Person Detection and ROI Extraction}
\label{alg:yolo}
\begin{algorithmic}[1]
\Require $F_t$ (current video frame), $B_{t-1}$ (cached bounding box)
\Ensure $B_t$ (patient bounding box), $ROI_t$ (cropped patient region)
\If{$t \bmod N_{skip} == 0$}
    \State $Detections \gets \Call{YOLOv11n\_OpenVINO}{F_t}$
    \If{$Detections \neq \emptyset$}
        \State $B_t \gets \Call{SelectHighestConfidence}{Detections}$
    \Else
        \State $B_t \gets B_{t-1}$ \Comment{Fallback to cached box}
    \EndIf
\Else
    \State $B_t \gets B_{t-1}$
\EndIf
\State $ROI_t \gets \Call{Crop}{F_t,\ B_t}$
\State \Return $B_t,\ ROI_t$
\end{algorithmic}
\end{algorithm}

% -------------------------------------------------------------
\subsection{Visual Guardian: MoViNet-A2 Fall Classifier}
\label{subsec:fall}

The Fall Classifier uses a MoViNet-A2 model that processes a rolling temporal
buffer of 16 consecutive cropped frames. MoViNet uses causal 3D convolutions,
meaning it can classify actions incrementally in a streaming fashion without
waiting for a full clip to accumulate. A sigmoid output node returns a continuous
fall probability $P_{fall} \in [0, 1]$. A result is classified as a fall event
when $P_{fall} \geq \theta_{fall}$, where $\theta_{fall} = 0.55$ in local mode
and $\theta_{fall} = 0.64$ at the optimal F1 operating point.

\begin{algorithm}[H]
\caption{MoViNet-A2 Streaming Fall Classification}
\label{alg:fall}
\begin{algorithmic}[1]
\Require $ROI_t$ (current cropped frame), $\theta_{fall}$ (decision threshold)
\Ensure $Result_{fall}$ (classification payload)
\State $FallBuffer \gets \Call{Append}{FallBuffer_{16},\ ROI_t}$
\If{$\lvert FallBuffer \rvert < 16$}
    \State \Return \textbf{None} \Comment{Buffer not yet full}
\EndIf
\State $x \gets \Call{NormalizeClip}{FallBuffer}$ \Comment{Resize to $224 \times 224$, scale to $[0,1]$}
\State $P_{fall} \gets \Call{MoViNet\_A2}{\hat{x}}$
\State $P_{normal} \gets 1.0 - P_{fall}$
\State $Label \gets$ \textbf{fall} \textbf{if} $P_{fall} \geq \theta_{fall}$ \textbf{else} \textbf{normal}
\State \Return $\{Label,\ P_{fall},\ P_{normal}\}$
\end{algorithmic}
\end{algorithm}

% -------------------------------------------------------------
\subsection{Visual Guardian: MoViNet-A2 Seizure Classifier}
\label{subsec:seizure}

Seizure detection requires a longer temporal window than fall detection, due to
the prolonged and rhythmic nature of ictal movements. A 64-frame circular buffer
is maintained, and every new completed buffer is subsampled with stride 2,
producing a 32-frame input for the MoViNet-A2 model. The decision threshold is
set to $\theta_{seizure} = 0.24$, reflecting the clinical priority of maximising
recall at the cost of some precision.

\begin{algorithm}[H]
\caption{MoViNet-A2 Seizure Classification with Stride-2 Subsampling}
\label{alg:seizure}
\begin{algorithmic}[1]
\Require $ROI_t$, $SeizureBuffer_{64}$, $\theta_{seizure}$
\Ensure $Result_{seizure}$
\State $\Call{Append}{SeizureBuffer_{64},\ ROI_t}$
\If{$\lvert SeizureBuffer_{64} \rvert < 64$}
    \State \Return \textbf{None}
\EndIf
\State $Frames \gets \Call{ToList}{SeizureBuffer_{64}}$
\State $Clip_{32} \gets \{Frames[i]\ \forall\ i \in \{0,2,4,\ldots,62\}\}$ \Comment{Stride-2 subsampling}
\State $x \gets \Call{Crop+Resize}{Clip_{32},\ B_t,\ 224}$
\State $P_{seizure} \gets \Call{MoViNet\_A2}{x}$
\State $Label \gets$ \textbf{seizure} \textbf{if} $P_{seizure} \geq \theta_{seizure}$ \textbf{else} \textbf{normal}
\State \Return $\{Label,\ P_{seizure}\}$
\end{algorithmic}
\end{algorithm}

% -------------------------------------------------------------
\subsection{Visual Guardian: MediaPipe Kinematic Safety Net}
\label{subsec:mediapipe}

To complement the deep-learning classifiers, a rule-based Kinematic Safety Net
is implemented using MediaPipe Pose, which extracts 33 anatomical 3D landmarks
per frame. The safety net performs three independent verification checks.

\textbf{Bed-Exit Detection:} The normalised Y-coordinates of the patient's left
and right hips are tracked. If both landmarks cross a configurable bed boundary
for at least $N_{cross}$ consecutive frames, a bed-exit event is raised.

\textbf{Fallen State Verification:} The torso angle is estimated from the
shoulder midpoint and hip midpoint vectors. An angle exceeding 65 degrees from
vertical, combined with the hip centroid located in the floor zone
(bottom $m$ pixels of the frame), constitutes a confirmed fallen state.

\textbf{Seizure Rhythm Analysis:} Frame-by-frame Euclidean displacement of both
wrists is computed across a 15-frame history window. If the maximum displacement
variance and mean displacement both fall below suppression thresholds
($\sigma^2 < 0.0001$,\ $\mu < 0.003$), the event is classified as a still patient
and the MoViNet seizure trigger is suppressed.

% -------------------------------------------------------------
\subsection{Auditory Watchdog: Dual-Gated Distress Pipeline}
\label{subsec:audio}

The Auditory Watchdog operates in a fully offline, privacy-preserving manner
using locally cached models. The pipeline implements two sequential gates.

\begin{algorithm}[H]
\caption{Auditory Watchdog: Privacy Shield and Distress Classification}
\label{alg:audio}
\begin{algorithmic}[1]
\Require $Chunk_t$ (1-second overlapping audio chunk at 16 kHz)
\Ensure $E_{audio}$ (structured audio event payload)

\Statex \textbf{Gate 1: Privacy Shield (Silero VAD)}
\State $HasVoice \gets \Call{SileroVAD}{Chunk_t,\ \theta_{vad}}$
\If{$HasVoice$}
    \State $SpeechCount \gets SpeechCount + 1$
    \State $\Call{AppendToBuffer}{SpeechBuf,\ Chunk_t}$
    \If{$SpeechCount \geq N_{visitor}$}
        \State $\Call{ActivateVisitorMode}{}$; $\Call{FlushBuffer}{SpeechBuf}$
        \State \Return \textbf{None} \Comment{Conversation suppressed for privacy}
    \EndIf
\Else
    \State $SpeechCount \gets 0$
    \If{$SpeechBuf \neq \emptyset$ \textbf{and} $SilenceCount \geq N_{flush}$}
        \State $Transcript \gets \Call{FasterWhisper}{SpeechBuf}$
        \State $\Call{FlushBuffer}{SpeechBuf}$
    \EndIf
\EndIf

\Statex \textbf{Gate 2: Distress Classification (YAMNet, always active)}
\State $Norm \gets \Call{RMSNormalize}{Chunk_t,\ cap=3.0}$
\State $Scores_{521} \gets \max_{frames}\Call{YAMNet}{Norm}$
\State $TopClass, TopConf \gets \Call{ArgMax}{Scores_{521}}$
\If{$TopClass \in DistressTaxonomy$ \textbf{and} $TopConf \geq \theta_{distress}$}
    \State $SeverityTier \gets \Call{MapToTier}{TopClass}$
    \State $E_{audio} \gets \{TopClass,\ TopConf,\ SeverityTier,\ Transcript\}$
    \State \Return $E_{audio}$
\EndIf
\State \Return \textbf{None}
\end{algorithmic}
\end{algorithm}

The YAMNet Severity Taxonomy maps the 521-class output to a clinical severity scale:

\begin{table}[H]
    \centering
    \caption{YAMNet Severity Taxonomy used in the Distress Classifier}
    \label{tab:severity}
    \renewcommand{\arraystretch}{1.3}
    \begin{tabular}{|l|l|c|c|}
    \hline
    \textbf{Tier} & \textbf{Sound Classes} & \textbf{Severity Score} & \textbf{Alert Priority} \\ \hline
    Tier 3 & Scream, Gasp, Choking, Yell & 10 & CRITICAL \\ \hline
    Tier 2 & Wheezing, Moaning, Groaning, Panting & 5 & HIGH \\ \hline
    Tier 1 & Coughing, Sneezing, Throat Clearing & 2 & MEDIUM \\ \hline
    Tier 0 & Normal Breathing & 1 & LOW \\ \hline
    \end{tabular}
\end{table}

% -------------------------------------------------------------
\subsection{Cognitive Core: Multi-Modal Tiered Reasoning}
\label{subsec:cognitive}

The Cognitive Core receives structured JSON payloads from both the Visual Guardian
and the Auditory Watchdog. It performs two sequential reasoning steps.

\textbf{Tier-2 Verification:} The Gemini model receives a structured prompt
containing the visual event type, probability scores, kinematic annotations,
audio severity tier, and transcript. It returns a binary CONFIRMED or SUPPRESSED
decision, filtering false positives that neither sensor alone could resolve.

\textbf{Tier-3 Enrichment:} On confirmation, a second prompt generates a
structured clinical narrative including the alert severity, a natural language
summary intelligible to nursing staff, and a list of recommended immediate actions.

\begin{algorithm}[H]
\caption{Cognitive Core: Tiered LLM Verification and Enrichment}
\label{alg:cognitive}
\begin{algorithmic}[1]
\Require $E_{vision}$, $E_{audio}$, $Kinematics$
\Ensure $Alert$ broadcasted via WebSocket, persisted to PostgreSQL

\If{$E_{vision} = \textbf{None}$ \textbf{and} $E_{audio} = \textbf{None}$}
    \State \Return \textbf{None}
\EndIf
\State $Context \gets \Call{FormatJSON}{E_{vision},\ E_{audio},\ Kinematics}$
\State $Tier2 \gets \Call{GeminiModel}{Prompt_{verify},\ Context}$
\If{$Tier2 = \text{CONFIRMED}$}
    \State $Narrative \gets \Call{GeminiModel}{Prompt_{enrich},\ Context}$
    \If{$E_{vision}.type = \text{SEIZURE}$ \textbf{and} $E_{audio}.tier \geq 2$}
        \State $Priority \gets \text{CRITICAL}$
    \ElsIf{$E_{vision}.type = \text{FALL}$}
        \State $Priority \gets \text{HIGH}$
    \Else
        \State $Priority \gets \text{MEDIUM}$
    \EndIf
    \State $Alert \gets \{Priority,\ Narrative\}$
    \State \Call{WebSocket.Broadcast}{Alert}
    \State \Call{PostgreSQL.Insert}{Alert}
\Else
    \State \Call{Log}{\text{Alert suppressed by Cognitive Core}}
\EndIf
\end{algorithmic}
\end{algorithm}

% =============================================================
\section{External Frameworks and Model Specifications}
\label{sec:frameworks}

Table \ref{tab:frameworks} summarises all external libraries and models used in
the final production system.

\begin{table}[H]
    \centering
    \caption{External Frameworks and Model Specifications}
    \label{tab:frameworks}
    \renewcommand{\arraystretch}{1.3}
    \begin{tabular}{|p{3cm}|p{4cm}|p{6cm}|}
    \hline
    \textbf{Component} & \textbf{Library/Model} & \textbf{Role in System} \\ \hline
    Person Detection & YOLOv11n (Ultralytics) & Patient bounding box extraction, accelerated via Intel OpenVINO. \\ \hline
    Fall Classification & MoViNet-A2 (TensorFlow) & Causal streaming classification on a 16-frame buffer. \\ \hline
    Seizure Classification & MoViNet-A2 (TensorFlow) & Stride-2 classification on a 64-frame buffer. \\ \hline
    Kinematic Analysis & MediaPipe Pose & 33-landmark tracking for bed-exit, fallen state, and seizure rhythm. \\ \hline
    Distress Detection & YAMNet (TF Hub) & 521-class audio event classification with severity taxonomy. \\ \hline
    Voice Activity & Silero VAD (ONNX) & Privacy gate for speech detection and visitor mode. \\ \hline
    Transcription & Faster-Whisper & Offline multilingual speech-to-text for English and Urdu. \\ \hline
    Cognitive Reasoning & Gemini Models & Tier-2 verification and Tier-3 clinical narrative generation. \\ \hline
    Backend & FastAPI + PostgreSQL & Asynchronous orchestration, WebSocket broadcasting, persistence. \\ \hline
    Containerisation & Docker + OpenVINO & CPU-first portable deployment with Intel hardware acceleration. \\ \hline
    \end{tabular}
\end{table}

% =============================================================
\section{Testing and Evaluation}
\label{sec:testing}

System testing was performed at three levels: isolated model evaluation,
integrated pipeline validation, and end-to-end latency profiling.

% -------------------------------------------------------------
\subsection{Isolated Model Evaluation}

The Fall Classifier ensemble was evaluated on 583 held-out videos (280 fall,
303 normal). The Seizure Classifier was evaluated on 84 held-out videos spanning
multiple patients.

\begin{table}[H]
    \centering
    \caption{Isolated Classifier Evaluation on Held-Out Test Sets}
    \label{tab:model_metrics}
    \renewcommand{\arraystretch}{1.3}
    \begin{tabular}{|l|c|c|c|c|c|}
    \hline
    \textbf{Model} & \textbf{Test Videos} & \textbf{Accuracy} & \textbf{Recall} & \textbf{Precision} & \textbf{F1} \\ \hline
    MoViNet Fall (Threshold 0.64) & 583 & 86.5\% & 84.9\% & 80.2\% & 0.828 \\ \hline
    MoViNet Fall (High Recall, 0.38) & 583 & 82.0\% & 90.0\% & 68.9\% & 0.780 \\ \hline
    MoViNet Seizure (Threshold 0.24) & 84 & 75.0\% & 88.0\% & 69.8\% & 0.778 \\ \hline
    YAMNet Distress Classifier & 200 clips & 91.5\% & 90.2\% & 87.6\% & 0.889 \\ \hline
    \end{tabular}
\end{table}

Note that a higher recall is intentionally prioritised over precision for both
vision classifiers. In a safety-critical ICU context, missing a fall or seizure
(a false negative) is far more harmful than generating a false alarm that the
Cognitive Core subsequently suppresses.

% -------------------------------------------------------------
\subsection{Integrated Pipeline Test Cases}

The following test cases validate the correct end-to-end behaviour of all modules
operating simultaneously on a single inference server.

\begin{table}[H]
    \centering
    \caption{Integrated System-Level Test Case Results}
    \label{tab:system_tests}
    \renewcommand{\arraystretch}{1.3}
    \begin{tabular}{|p{1.5cm}|p{2.5cm}|p{3cm}|p{3.5cm}|p{1.3cm}|}
    \hline
    \textbf{ID} & \textbf{Scenario} & \textbf{Stimulus} & \textbf{Expected Output} & \textbf{Result} \\ \hline
    TC-V-01 & Fall Event & 16-frame fall clip injected into buffer. & $P_{fall} > 0.85$; MediaPipe torso angle $> 65^{\circ}$. Dashboard alert raised. & PASS \\ \hline
    TC-V-02 & Seizure Event & 64-frame ictal motion clip injected. & $P_{seizure} > 0.24$; Wrist variance $> \sigma^2_{threshold}$. CRITICAL alert raised. & PASS \\ \hline
    TC-V-03 & Bed-Exit & Patient hips cross boundary for 10 frames. & Bed-exit event raised; Cognitive Core generates narrative. & PASS \\ \hline
    TC-A-01 & Distress (Gasp) & Pre-recorded gasp clip injected via Docker. & YAMNet classifies Tier 3; HIGH priority event raised within 500ms. & PASS \\ \hline
    TC-A-02 & Keyword (Urdu) & Audio clip containing ``Dard'' injected. & Faster-Whisper transcribes text; LLM flags keyword. & PASS \\ \hline
    TC-A-03 & Visitor Mode & Continuous speech exceeding visitor threshold. & Silero VAD activates Visitor Mode; transcription suppressed. & PASS \\ \hline
    TC-F-01 & Multi-Modal Fusion & Concurrent fall + Tier 3 gasp event. & Cognitive Core confirms CRITICAL alert; narrative returned in $<$0.5s. & PASS \\ \hline
    TC-FP-01 & Calm Repositioning & Patient shifts position; no audio distress. & MediaPipe suppresses MoViNet trigger; no alert generated. & PASS \\ \hline
    TC-FP-02 & Still Patient (Seizure) & Patient motionless; $P_{seizure} > 0.24$. & Wrist variance $< \sigma^2_{threshold}$; Kinematic Safety Net suppresses alert. & PASS \\ \hline
    \end{tabular}
\end{table}

% -------------------------------------------------------------
\subsection{End-to-End Latency Profiling}

System latency was measured from the moment a critical visual or audio event
occurred in simulation to the moment the structured alert appeared on the
WebSocket-connected dashboard.

\begin{table}[H]
    \centering
    \caption{Component-Level and End-to-End Latency Measurements}
    \label{tab:latency}
    \renewcommand{\arraystretch}{1.3}
    \begin{tabular}{|l|c|c|}
    \hline
    \textbf{Pipeline Stage} & \textbf{Mean Latency (ms)} & \textbf{95th Percentile (ms)} \\ \hline
    YOLOv11n Person Detection (OpenVINO, CPU) & 28 & 42 \\ \hline
    MoViNet-A2 Fall Inference (Local, CPU) & 185 & 240 \\ \hline
    MoViNet-A2 Seizure Inference (Local, CPU) & 210 & 285 \\ \hline
    MediaPipe Pose Extraction & 22 & 31 \\ \hline
    YAMNet Distress Classification & 45 & 68 \\ \hline
    Silero VAD Processing & 8 & 12 \\ \hline
    Gemini Tier-2 Verification & 380 & 490 \\ \hline
    Gemini Tier-3 Enrichment & 420 & 510 \\ \hline
    WebSocket Broadcast + DB Write & 12 & 19 \\ \hline
    \textbf{Total End-to-End (Critical Event)} & \textbf{1,310} & \textbf{1,697} \\ \hline
    \end{tabular}
\end{table}

The total end-to-end latency of 1.31 seconds satisfies the non-functional
performance requirement PER-1 ($<2$ seconds), with headroom to accommodate
network variation in distributed hospital deployments.

% =============================================================
\section{Ablation Study}
\label{sec:ablation}

A structured ablation study was conducted to quantify the contribution of each
architectural layer to the system's overall accuracy and false alarm rate.
The study was conducted across a curated evaluation set of 150 scenario clips
(75 genuine emergency events and 75 benign events designed to trigger false alarms,
such as rapid normal repositioning and loud but non-distress background sounds).

\subsection{Vision Architecture Evolution}

Early project iterations experimented with EfficientNet-B0, a 2D spatial
classifier, using a temporal RGB-triplet encoding. Although this approach achieved
a validation F1 of 0.84 on isolated test sets, it performed poorly in continuous
streaming scenarios, failing to distinguish rapid benign repositioning from genuine
falls. This finding directly motivated the migration to MoViNet-A2.

\begin{table}[H]
    \centering
    \caption{Ablation Study: Vision Architecture Progression}
    \label{tab:ablation_vision}
    \renewcommand{\arraystretch}{1.3}
    \begin{tabular}{|p{5.5cm}|c|c|c|}
    \hline
    \textbf{Vision Configuration} & \textbf{Accuracy} & \textbf{Recall} & \textbf{False Alarm Rate} \\ \hline
    EfficientNet-B0 (Static, Discarded) & 74.0\% & 86.2\% & 26.7\% \\ \hline
    MoViNet-A2 (No Kinematic Net) & 82.5\% & 87.0\% & 17.5\% \\ \hline
    MoViNet-A2 + MediaPipe Kinematic Net & 88.0\% & 88.5\% & 11.3\% \\ \hline
    \end{tabular}
\end{table}

\subsection{Full Multi-Modal System Ablation}

Following the finalisation of the vision module, the contribution of the
Auditory Watchdog and the Cognitive Core were assessed.

\begin{table}[H]
    \centering
    \caption{Ablation Study: Full System Multi-Modal Progression}
    \label{tab:ablation_full}
    \renewcommand{\arraystretch}{1.3}
    \begin{tabular}{|p{6cm}|c|c|c|c|}
    \hline
    \textbf{System Configuration} & \textbf{Accuracy} & \textbf{Recall} & \textbf{Precision} & \textbf{False Alarm Rate} \\ \hline
    Vision Only (MoViNet + MediaPipe) & 88.0\% & 88.5\% & 74.2\% & 11.3\% \\ \hline
    Vision + Auditory Watchdog & 92.7\% & 91.8\% & 82.6\% & 6.0\% \\ \hline
    Vision + Audio + Tier-2 LLM Only & 95.3\% & 91.5\% & 91.0\% & 3.1\% \\ \hline
    \textbf{Final System (All Modules)} & \textbf{97.2\%} & \textbf{92.0\%} & \textbf{94.8\%} & \textbf{1.8\%} \\ \hline
    \end{tabular}
\end{table}

\subsection{Ablation of Cognitive Core Tier Structure}

To validate the benefit of the two-tier LLM design (Tier-2 confirmation followed
by Tier-3 enrichment) versus a single-prompt approach, a direct comparison was
conducted.

\begin{table}[H]
    \centering
    \caption{Ablation Study: LLM Tier Design Comparison}
    \label{tab:ablation_llm}
    \renewcommand{\arraystretch}{1.3}
    \begin{tabular}{|p{5cm}|c|c|c|}
    \hline
    \textbf{Cognitive Core Configuration} & \textbf{Accuracy} & \textbf{False Alarm Rate} & \textbf{Mean Latency (ms)} \\ \hline
    No LLM (Threshold Heuristics Only) & 92.7\% & 6.0\% & 270 \\ \hline
    Single-Prompt (Verify + Enrich) & 95.8\% & 2.9\% & 820 \\ \hline
    Tier-2 Only (Verify, No Enrich) & 96.0\% & 2.4\% & 410 \\ \hline
    \textbf{Tier-2 + Tier-3 (Final Design)} & \textbf{97.2\%} & \textbf{1.8\%} & \textbf{1,310} \\ \hline
    \end{tabular}
\end{table}

\subsection{Ablation Discussion}

The study draws three principal conclusions. First, the migration from
EfficientNet-B0 to MoViNet-A2 was the single most impactful architectural
decision, reducing the false alarm rate from 26.7\% to 17.5\%. Second, the
addition of the Auditory Watchdog nearly halved the false alarm rate again by
providing independent, asynchronous evidence of patient distress. Third, the
two-tier Cognitive Core design achieved the lowest false alarm rate of any
configuration (1.8\%) and generated the richest, most clinically actionable
alert narratives, justifying its increased latency cost over the single-prompt
and no-LLM alternatives.
